% Part: model-theory
% Chapter: basics
% Section: isomorphism

\documentclass[../../../include/open-logic-section]{subfiles}

\begin{document}

\olfileid{mod}{bas}{iso}
\olsection{એકરૂપ સંરચનાઓ}

પ્રથમ-ક્રમ !!{structure}s બે રીતે એકસરખી હોઈ શકે છે. એક રીતે તેઓ
એકસરખી હોય છે જો તેઓ એ જ !!{sentence}sને સત્ય બનાવે. આવી
!!{structure}sને આપણે \emph{પ્રાથમિક રીતે સમકક્ષ} કહીએ છીએ. પરંતુ બહુ
જુદી સંરચનાઓ પણ એ જ !!{sentence}sને સત્ય બનાવી શકે છે---દાખલા તરીકે,
એક !!{enumerable} હોઈ શકે અને બીજી નહિ. તેનું કારણ એ છે કે
!!a{structure}નાં ઘણાં લક્ષણો પ્રથમ-ક્રમ ભાષાઓમાં વ્યક્ત કરી શકાતાં નથી;
કાં તો ભાષા પૂરતી સમૃદ્ધ નથી અથવા લેવેનહાઇમ--સ્કોલેમ પ્રમેય જેવી
પ્રથમ-ક્રમ તર્કશાસ્ત્રની મૂળભૂત મર્યાદાઓ તેને રોકે છે. તેથી
!!{structure}s વધુ ચુસ્ત અર્થમાં પણ એકસરખી હોઈ શકે છે: તેમને બનાવતી
વસ્તુઓ જ જુદી હોય, પણ તેમનાં સંરચનાત્મક લક્ષણો જુદાં ન હોય. આ વિચારને
ચોક્કસ બનાવવાની એક રીત \emph{એકરૂપતા}નો ખ્યાલ છે.

\begin{defn}
\ollabel{defn:elem-equiv}
એક જ !!{language}~$\Lang{L}$ માટેની બે !!{structure}s $\Struct{M}$ અને
$\Struct M'$ આપેલી હોય ત્યારે, $\Lang{L}$ના દરેક !!{sentence}~$!A$ માટે
$\Sat{M}{!A}$ ત્યારે અને માત્ર ત્યારે જ્યારે $\Sat{M'}{!A}$ હોય, તો અને
માત્ર તો $\Struct{M}$ને $\Struct M'$ની \emph{પ્રાથમિક રીતે સમકક્ષ}
કહીએ છીએ અને $\Struct{M} \equiv \Struct M'$ લખીએ છીએ.
\end{defn}

\begin{defn}
\ollabel{defn:isomorphism} એક જ !!{language}~$\Lang L$ માટેની બે
!!{structure}s $\Struct{M}$ અને $\Struct M'$ આપેલી હોય ત્યારે, નીચેની
શરતો સંતોષતું વિધેય $h \colon \Domain{M} \to \Domain{M'}$ હોય તો અને
માત્ર તો $\Struct{M}$ને~$\Struct M'$ની \emph{એકરૂપ} કહીએ છીએ અને
$\Struct{M} \simeq \Struct M'$ લખીએ છીએ:
\begin{enumerate}
\item $h$ !!{injective} છે: જો $h(x) = h(y)$, તો $x = y$;
\item $h$ !!{surjective} છે: દરેક $y \in \Domain{M'}$ માટે એવો
  $x \in \Domain{M}$ છે કે $h(x) = y$;
\item \ollabel{defn:iso-const}દરેક !!{constant} $c$ માટે:
  $h(\Assign{c}{M}) = \Assign{c}{M'}$;
\item \ollabel{defn:iso-pred}દરેક $n$-સ્થાની !!{predicate}~$P$ માટે:
  \[
  \tuple{a_1, \dots, a_n}\in \Assign{P}{M} \quad\text{ત્યારે અને માત્ર ત્યારે}\quad
  \tuple{h(a_1), \dots, h(a_n)} \in \Assign{P}{M'};
  \]
\item \ollabel{defn:iso-func}દરેક $n$-સ્થાની !!{function} $f$ માટે:
  \[
  h(\Assign{f}{M}(a_1, \dots, a_n)) =
  \Assign{f}{M'}(h(a_1), \dots, h(a_n)).
  \]
\end{enumerate}
\end{defn}

\begin{thm}
\ollabel{thm:isom}
જો $\Struct{M} \iso \Struct M'$, તો $\Struct{M} \elemequiv
\Struct{M'}$.
\end{thm}

\begin{proof}
ધારો કે $h$ એ $\Struct{M}$થી $\Struct M'$ પરની એકરૂપતા છે. કોઈ પણ
નિયુક્તિ~$s$ માટે $h \circ s$ એ $h$ અને $s$નું સંયોજન છે, એટલે કે
$\Struct{M'}$માં એવી નિયુક્તિ કે $(h \circ s)(x) = h(s(x))$.
$t$ અને $!A$ પરના આગમનથી નીચેના વધુ પ્રબળ દાવા સાબિત કરી શકાય છે:
\begin{enumerate}
  \item[a.] $h(\Value{t}{M}[s]) = \Value{t}{M'}[h\circ s]$.
  \item[b.] $\Sat{M}{!A}[s]$ ત્યારે અને માત્ર ત્યારે જ્યારે
    $\Sat{M'}{!A}[h \circ s]$.
\end{enumerate}
પહેલો દાવો~$t$ની જટિલતા પરના આગમનથી સાબિત થાય છે.
\begin{enumerate}
\item જો $t \ident c$, તો $\Value{c}{M}[s] = \Assign{c}{M}$ અને
  $\Value{c}{M'}[h \circ s] = \Assign{c}{M'}$. તેથી
  $h(\Value{t}{M}[s]) = h(\Assign{c}{M}) = \Assign{c}{M'}$
  (\olref{defn:isomorphism}ની \olref{defn:iso-const} મુજબ) $=
  \Value{t}{M'}[h \circ s]$.
\item જો $t \ident x$, તો $\Value{x}{M}[s] = s(x)$ અને
  $\Value{x}{M'}[h \circ s] = h(s(x))$. તેથી
  $h(\Value{x}{M}[s]) = h(s(x)) = \Value{x}{M'}[h \circ s]$.
\item જો $t \ident f(t_1, \dots, t_n)$, તો
  \begin{align*}
    \Value{t}{M}[s] & = \Assign{f}{M}(\Value{t_1}{M}[s], \dots, \Value{t_n}{M}[s]) \quad\text{અને}\\
  \Value{t}{M'}[h \circ s] & = \Assign{f}{M'}(\Value{t_1}{M'}[h \circ
    s], \dots, \Value{t_n}{M'}[h \circ s]).
  \end{align*}
  આગમન પરિકલ્પના મુજબ દરેક $i$ માટે $h(\Value{t_i}{M}[s])
  = \Value{t_i}{M'}[h\circ s]$. તેથી
  \begin{align}
    h(\Value{t}{M}[s])
    & = h(\Assign{f}{M}(\Value{t_1}{M}[s], \dots, \Value{t_n}{M}[s])) \notag\\
    & = \Assign{f}{M'}(h(\Value{t_1}{M}[s]), \dots,
    h(\Value{t_n}{M}[s])) \ollabel{iso-1}\\
    & = \Assign{f}{M'}(\Value{t_1}{M'}[h \circ s], \dots,
    \Value{t_n}{M'}[h \circ s]) \ollabel{iso-2}\\
    & = \Value{t}{M'}[h\circ s] \notag
  \end{align}
  અહીં \olref{iso-1}, \olref{defn:isomorphism}ની
  \olref{defn:iso-func}માંથી અને \olref{iso-2} આગમન પરિકલ્પનામાંથી ફલિત થાય છે.
  \sourcecorrection{OLMTB-003}{સ્થિર અંગ્રેજી મૂળની
    \texttt{isomorphism.tex}, પંક્તિઓ~89--90માં
    $\Value{t}{M'}[h\circ s]$ની વ્યાખ્યામાં વિધેયનું અર્થઘટન ભૂલથી
    $\Assign{f}{M}$ લખાયું છે. અહીં લક્ષ્ય સંરચનાનું
    $\Assign{f}{M'}$ પુનઃસ્થાપિત કર્યું છે.}
  \sourcecorrection{OLMTB-004}{એ જ મૂળની પંક્તિ~96માં
    $h(\Assign{f}{M}(\dots))$નું બહારનું જમણું કૌંસ ગાયબ છે. અહીં તે
    કૌંસ પુનઃસ્થાપિત કર્યું છે.}
\end{enumerate}
ભાગ~(b)ની સાબિતી અભ્યાસપ્રશ્ન તરીકે છોડી છે.

જો $!A$ વાક્ય હોય, તો નિયુક્તિઓ~$s$ અને $h \circ s$ અપ્રસ્તુત છે અને
આપણે $\Sat{M}{!A}$ ત્યારે અને માત્ર ત્યારે જ્યારે $\Sat{M'}{!A}$ મેળવીએ છીએ.
\end{proof}

\begin{prob}
\olref[mod][bas][iso]{thm:isom}ના ભાગ~(b)ની સાબિતી વિગતે પૂરી કરો.
\olref[mod][bas][iso]{defn:isomorphism}માં એકરૂપતાઓને લક્ષણિત કરતા પાંચ
ગુણધર્મોમાંથી દરેકનો ઉપયોગ ક્યાં થાય છે તે ખાસ નોંધો.
\end{prob}

\begin{defn}
સંરચના $\Struct{M}$ની પોતાના પરની એકરૂપતાને $\Struct{M}$ની
\emph{સ્વ-એકરૂપતા} કહીએ છીએ.
\end{defn}

\begin{prob}
બતાવો કે કોઈ પણ સંરચના $\Struct{M}$ માટે, જો $X$ એ $\Struct{M}$નો
વ્યાખ્યેય ઉપગણ હોય અને $h$ એ $\Struct{M}$ની સ્વ-એકરૂપતા હોય, તો
$X = \Setabs{h(x)}{x \in X}$ (એટલે કે $X$, $h$ હેઠળ સ્થિર છે).
\end{prob}


\end{document}
